\scanpage{90}
\mspage{9}
% archive foliation: 89

action tautologique de $\unsure{\mathfrak{S}|\mathcal G}_3$, et sur $\unsure{\boldsymbol{\mu}|\mu}_4^*$ via le caractère mixte $\operatorname{sg}(g)\cdot\unsure{\zeta}^2$ $(g\in\unsure{\mathfrak{S}|\mathcal G}_3,\unsure{\zeta}\in\unsure{\boldsymbol{\mu}|\mu}_4)$. Donc le sous-groupe de $\unsure{\mathfrak{S}|\mathcal G}_3\times\unsure{\boldsymbol{\mu}|\mu}_4$ qui opère trivialement sur $\unsure{\boldsymbol{\mu}|\mu}_4^*$ est
\[
(48)\qquad\mathbb{D}'_{3S_0}=(\unsure{\mathfrak{S}|\mathcal G}_3)_{S_0}\times_{\unsure{\boldsymbol{\mu}|\mu}_{2S_0}}(\unsure{\boldsymbol{\mu}|\mu}_4)_{S_0}.
\]
Si on prend comme base absolue non $S_0$, mais $\widetilde{S}_0=(\unsure{\boldsymbol{\mu}|\mu}_4^*)_{S_0}$ de sorte que l'on trouve, sur les schémas\inserted{/$\widetilde{S}_0$} envisagés
\[
\left\{\begin{array}{l}
\unsure{\boldsymbol{\mu}|\mu}_{4S}^*\deleted{\unsure{S}}\text{ muni d'une section canonique }i_S\\[4pt]
\unsure{\boldsymbol{\mu}|\mu}_{4S}\simeq(\mathbb{Z}/4\mathbb{Z})_S
\end{array}\right.
\]
alors le groupe (48) s'identifie au groupe constant sur $S$ \deleted{\unsure{de valeur}}\inserted{défini par le groupe \unsure{abstrait}}
\[
(49)\qquad\mathbb{D}'_3=\unsure{\mathfrak{S}|\mathcal G}_3\times_{\mathbb{Z}/2\mathbb{Z}}(\mathbb{Z}/4\mathbb{Z})
\]
qui est une extension de $\mathbb{D}_3\simeq\unsure{\mathfrak{S}|\mathcal G}_3$ par $\mathbb{Z}/2\mathbb{Z}\simeq\{\pm1\}$
\[
(50)\qquad\diagram{\begin{tikzcd}[column sep=small]
1\arrow[r,no head]&\deleted{\unsure{\mathbb{Z}/2\mathbb{Z}}}\{\pm1\}\arrow[r,no head]&\mathbb{D}'_3\arrow[r,no head]&\mathbb{D}_3\arrow[r,no head]&1
\end{tikzcd}}.
\]
\deleted{\unsure{Voyons}\ \unsure{si}}
\deleted{L'action envisagée de $\unsure{\mathfrak{S}|\mathcal G}_3\times\unsure{\boldsymbol{\mu}|\mu}_4$}\inserted{Dans le \unsure{cas présent} i.e. en} travaillant sur $\widetilde{S}$ i.e. avec un choix fixé de $i=\sqrt{-1}$, on peut distinguer
